\documentclass[11pt]{article}
\usepackage[T1]{fontenc}
\usepackage[margin=1in]{geometry}
\usepackage{amsmath,amssymb,amsthm,mathtools}
\usepackage{graphicx}
\usepackage[hidelinks]{hyperref}
\usepackage{microtype}
\usepackage{xcolor}

\newtheorem{theorem}{Theorem}
\newtheorem{remark}{Remark}
\newcommand{\M}{\mathcal M}
\newcommand{\C}{\mathbb C}
\newcommand{\tr}{\operatorname{tr}}
\newcommand{\diag}{\operatorname{diag}}
\newcommand{\rank}{\operatorname{rank}}

\title{Schatten--\(p\) approximation selects the strict spectral
approximant for every matrix hyperplane}
\author{Candidate partial result for arXiv:2603.07498}
\date{9 August 2026}

\begin{document}
\maketitle

\noindent\textbf{Status.}
\texttt{candidate\_partial\_likely\_valid}.  The source conjecture is not
settled for arbitrary subspaces.  We prove it for every complex
codimension-one subspace, in all finite rectangular dimensions.

\section{Source problem and result}

For \(A\in M_{m\times n}(\C)\) and a complex linear subspace \(\M\), let
\(Y_p\in\M\) minimize the Schatten \(p\)-norm \(\|A-Y\|_p\), where
\(1<p<\infty\).  Grover and Gupta ask in their Conjecture~(5) whether
\[
                  Y_p\longrightarrow Y^{(\mathrm{st})}
                  \qquad (p\to\infty),
\]
where \(Y^{(\mathrm{st})}\) is the strict spectral approximant: among
spectral-norm minimizers, it lexicographically minimizes the complete
singular-value vector of the residual.

\begin{center}
\includegraphics[width=.62\textwidth]{figures/open_problem_crop.png}

\small Real crop from page 2 of the source paper, including Conjecture~(5).
\end{center}

\paragraph{Proof intuition.}
A hyperplane has a single normal matrix \(G\).  Schatten duality makes the
finite-\(p\) residual the unique equality case in one H\"older inequality,
so it can be written explicitly from the singular values of \(G\).  At
\(p=\infty\), equality in nuclear/operator duality forces the same maximal
singular block in \emph{every} spectral minimizer.  The strict minimizer is
then obtained by setting the unconstrained complementary block equal to
zero.  The explicit finite-\(p\) formula converges to precisely that
matrix.

\section{The hyperplane theorem}

\begin{theorem}
Let \(G\in M_{m\times n}(\C)\setminus\{0\}\) and
\[
       \M=\{X\in M_{m\times n}(\C):\tr(G^*X)=0\}.
\]
For \(A\in M_{m\times n}(\C)\), let \(Y_p\) be its unique Schatten
\(p\)-best approximant from \(\M\), \(1<p<\infty\).  Then
\[
                         Y_p\longrightarrow Y^{(\mathrm{st})}.
\]
More explicitly, set \(c=\tr(G^*A)\).  If \(c=0\), then
\(Y_p=Y^{(\mathrm{st})}=A\).  If \(c\ne0\), take a thin singular value
decomposition
\[
 G=U_r\diag(g_1,\ldots,g_r)V_r^*,\qquad
 g_i>0,\quad r=\rank G,
\]
and set \(\omega=c/|c|\), \(q=p/(p-1)\).  Then
\begin{align}
 A-Y_p
   &=\frac{\omega |c|}{\sum_{i=1}^r g_i^q}\,
     U_r\diag\!\left(g_1^{q-1},\ldots,g_r^{q-1}\right)V_r^*,       \tag{1}\\
 A-Y^{(\mathrm{st})}
   &=\frac{\omega |c|}{\sum_{i=1}^r g_i}\,U_rV_r^*.              \tag{2}
\end{align}
\end{theorem}

\begin{proof}
Approximating \(A\) from \(\M\) is equivalent to minimizing
\(\|R\|_p\) subject to
\[
                         \tr(G^*R)=c,                             \tag{3}
\]
where \(R=A-Y\).  The assertion for \(c=0\) is immediate, so suppose
\(c\ne0\).

Schatten H\"older duality, with \(q=p/(p-1)\), gives every feasible
\(R\) the lower bound
\[
 |c|=|\tr(G^*R)|\leq \|G\|_q\|R\|_p,\qquad
 \|G\|_q=\Bigl(\sum_i g_i^q\Bigr)^{1/q}.                          \tag{4}
\]
The matrix on the right side of (1) is feasible and attains equality
in (4), hence is a minimizer.  Since the Schatten \(p\)-norm is strictly
convex for \(1<p<\infty\), the minimizer is unique.  This proves (1).

It remains to identify the strict spectral residual.  Put
\[
                         \rho=\frac{|c|}{\|G\|_1}
                              =\frac{|c|}{\sum_i g_i}.
\]
Nuclear/operator duality gives \(\|R\|_\infty\geq\rho\) for every
feasible \(R\), while
\[
                         R_0=\omega\rho\,U_rV_r^*                 \tag{5}
\]
is feasible and has operator norm \(\rho\).  Thus \(\rho\) is the
optimal spectral error.

We now determine all equality cases.  Write \(u_i,v_i\) for the
singular vectors in the displayed decomposition of \(G\).  If \(R\) is
any spectral minimizer, then
\begin{equation}
\begin{split}
 |c|
 &=\left|\sum_{i=1}^r g_i\,u_i^*Rv_i\right|\\
 &\leq\sum_{i=1}^r g_i|u_i^*Rv_i|
 \leq\rho\sum_{i=1}^r g_i
 =|c|.
\end{split}
\tag{6}
\end{equation}
All inequalities are equalities.  Because every \(g_i>0\), equality in
the last inequality yields \(|u_i^*Rv_i|=\rho\) for every \(i\), and
equality in the triangle inequality gives their common phase:
\[
                         u_i^*Rv_i=\omega\rho.                    \tag{7}
\]
Moreover,
\[
 \rho=|u_i^*Rv_i|\leq\|Rv_i\|_2\leq\|R\|_\infty=\rho.
\]
Equality in Cauchy--Schwarz implies \(Rv_i=\omega\rho u_i\).
Applying the same argument to \(R^*\) gives
\(R^*u_i=\overline\omega\rho v_i\).  Consequently, relative to the
orthogonal decompositions
\[
 \C^n=\operatorname{span}\{v_i\}_{i=1}^r\oplus V^\perp,\qquad
 \C^m=\operatorname{span}\{u_i\}_{i=1}^r\oplus U^\perp,
\]
every spectral-minimizing residual has block form
\[
                         R=\begin{pmatrix}
                              \omega\rho I_r&0\\[2pt]0&B
                            \end{pmatrix},
                    \qquad \|B\|_\infty\leq\rho.                  \tag{8}
\]
Its singular-value vector therefore begins with at least \(r\) copies
of \(\rho\).  Formula (5) corresponds to \(B=0\), so its vector is
\[
                 (\underbrace{\rho,\ldots,\rho}_{r\ \mathrm{times}},
                  0,\ldots,0),
\]
which is lexicographically smallest among (8).  It is also the unique
matrix with this vector, because zero remaining singular values force
\(B=0\).  Hence \(R_0=A-Y^{(\mathrm{st})}\), proving (2).

Finally, \(q\downarrow1\) as \(p\to\infty\), so
\(g_i^{q-1}\to1\) and \(\sum_i g_i^q\to\sum_i g_i\).  Formula (1)
therefore converges to (2).
\end{proof}

\section{Why this is beyond the unique spectral case}

Rank-deficient \(G\) produces genuinely nonunique spectral
approximations.  For example, take
\[
 G=\diag(2,1,0),\qquad A=\diag(1,1,0),\qquad
 \M=\ker\!\bigl(X\mapsto\tr(G^*X)\bigr).
\]
Here \(c=3\), \(\rho=1\), and all residuals
\(\diag(1,1,z)\), \(|z|\leq1\), are spectral minimizers.  The strict
residual is \(\diag(1,1,0)\), whereas (1) gives
\[
 A-Y_p=\frac{3}{2^q+1}\diag(2^{q-1},1,0)
          \longrightarrow\diag(1,1,0).
\]

\section{Verification, limitations, and novelty}

\paragraph{Computational check.}
The accompanying script generated \(1{,}000\) random complex
hyperplanes of sizes \(3\times4\) and \(4\times4\), including all
possible ranks.  At \(p=2,4,10,50,200\), it checked feasibility,
compared (1) against random feasible perturbations, and checked
convergence to (2).  The maximum feasibility error was
\(5.024\cdot10^{-15}\), the minimum sampled optimality slack was
\(5.813\cdot10^{-7}\), and the maximum \(p=200\) Frobenius error to the
strict residual was \(3.231\cdot10^{-2}\).  These tests are not used as
proof.

\paragraph{Limitations.}
The unrestricted conjecture remains open.  In codimension at least two,
the dual norming matrix can vary with \(q\), so the single-normal
equality argument above does not supply a uniform lexicographic
selection theorem.

\paragraph{Novelty check.}
A bounded search of the run indexes, the source paper, and directly
related older literature found no statement of this all-hyperplane
theorem.  Ziętak's pointwise eventual comparison theorem for a fixed
competitor does not by itself control the moving minimizers \(Y_p\).
This is a preliminary novelty assessment, not a literature claim.

\paragraph{Human-review recommendation.}
Check especially the equality passage from (6) to the block form (8)
and search for an existing codimension-one theorem stated in duality
language.  If both checks pass, the result appears ready to promote as
a rigorous partial resolution.

\begin{thebibliography}{9}
\bibitem{GG}
P.~Grover and K.~K.~Gupta,
\emph{Properties of best approximations with respect to the Ky Fan
\(p\)-\(k\) norm, and the strict spectral approximant of a matrix},
arXiv:2603.07498v2 (2026).

\bibitem{Zietak}
K.~Ziętak,
\emph{Strict spectral approximation of a matrix and some related
problems},
Applicationes Mathematicae \textbf{44} (2017), 267--280.
\end{thebibliography}

\end{document}
